\section{Curvature}

\subsection{Families of curves}

\bdefn

    Given intervals $I,J\subseteq\bR$, a continuous map $\Gamma\colon J\times I\to M$ is called a {\emphcolor one-parameter
    family of curves}.
    This defines two collections of curves:
    \benum
        \item the {\emphcolor main curves}: for each $s\in J$ the curve $\Gamma_s\colon I\to M$ by $\Gamma_s(t)=\Gamma(s,t)$.
        \item the {\emphcolor transverse curves}: for $t\in I$, the curve $\Gamma^{(t)}\colon J\to M$ by
        $\Gamma^{(t)}(s)=\Gamma(s,t)$.
    \eenum

\edefn

If $\Gamma$ is smooth, then denote the velocity vectors of the main and transverse curves by
$$ \partial_t\Gamma(s,t) = (\Gamma_s)'(t),\qquad \partial_s\Gamma(s,t) = (\Gamma^{(t)})'(s) , $$
both of which lie in $T_{\Gamma(s,t)}M$.
These are vector fields along $\Gamma$ (smooth when $\Gamma$ is smooth).

\bdefn

    A curve $\gamma\colon I\to M$ is {\emphcolor regular} if $\gamma'$ is nonzero.
    A curve $\gamma\colon[a,b]\to M$ is {\emphcolor admissible} if there is a partition $a_0<\cdots<a_k$ of $[a,b]$
    such that $\gamma|_{[a_i,a_{i+1}]}$ is regular for all $i$.

    A one-parameter family of curves $\Gamma$ is called {\emphcolor admissible} if its domain is of the form
    $J\times[a,b]$ for some open interval $J$, and there is a partition $a_0<\cdots<a_k$ of $[a,b]$ such that
    $\Gamma$ is smooth on each rectangle $J\times[a_i,a_{i+1}]$, and finally $\Gamma_s$ is an admissible curve.

\edefn

For an admissible family $\Gamma$, the transverse curves $\Gamma^{(t)}$ are smooth (as $\Gamma$ is smooth on
$J\times\set t$), but the main curves are only piecewise regular.
Thus $\partial_s\Gamma,\partial_t\Gamma$ are smooth on the rectangles $J\times[a_i,a_{i+1}]$ but not necessarily on the whole
domain.

If $\Gamma$ is an admissible family, then $\partial_s\Gamma$ is a piecewise smooth vector field along $\Gamma$.
This means it is continuous, and its restriction to each $J\times[a_i,a_{i+1}]$ is smooth.
As explained its restrictions are smooth, and we explain why it is continuous.

Note that for each $i$, $\partial_s\Gamma$'s values on $J\times\set{a_i}$ depend only on the values of $\Gamma$
on that set, since the derivative is taken with respect to $s$.
And $\partial_s\Gamma$ is smooth on $J\times[a_{i-1},a_i]$ and $J\times[a_i,a_{i+1}]$, and their values must agree on their
overlap (taking the limit as $t$ goes to $a_i$).

Conversely, $\partial_t\Gamma$ is not in general continuous at $t=a_i$.

\bdefn

    If $\gamma\colon[a,b]\to M$ is an admissible curve, a {\emphcolor variation} of $\gamma$ is an admissible family
    of curves $\Gamma\colon J\times[a,b]\to M$ such that $J$ is an open interval containing $0$ and $\Gamma_0=\gamma$.

    It is a {\emphcolor proper variation} if in addition all the main curves have the same endpoints:
    $\Gamma_s(a)=\gamma(a)$ and $\Gamma_s(b)=\gamma(b)$ for all $s\in J$.

    If $\Gamma$ is a variation of $\gamma$, then the {\emphcolor variation field} of $\Gamma$ is the smooth vector
    field $V(t)=\partial_s\Gamma(0,t)$ along $\gamma$.
    We say that a vector field $V$ along $\gamma$ is {\emphcolor proper} if $V(a)$ and $V(b)$ are both zero.

\edefn

\blemm

    If $\gamma$ is an admissible curve and $V$ a piecewise smooth vector field along $\gamma$, then $V$ is the
    variation field of some variation of $\gamma$.
    If $V$ is proper, the variation can be taken to be proper as well.

\elemm

\bproof

    Let $\Gamma(s,t)=\exp_{\gamma(t)}(sV(t))$.
    By compactness of $[a,b]$ there is some $\epsilon>0$ such that $\Gamma$ is defined on
    $(-\epsilon,\epsilon)\times[a,b]$.
    Then $\Gamma$ is smooth on each $(-\epsilon,\epsilon)\times[a_i,a_{i+1}]$ when $V$ is smooth on $[a_i,a_{i+1}]$,
    and it is continuous on its whole domain.

    The variation field of $\Gamma$ is $V$:
    $$ \partial_s\Gamma(0,t) = \partial_s|_{s=0}\exp_{\gamma(t)}(sV(t)) = V(t) $$
    because $\exp_p(sv)$ is the geodesic with initial velocity $v$.

    Furthermore if $V$ is proper, so $V(a),V(b)$ are zero, then
    $$ \Gamma(s,a) = \exp_{\gamma(a)}(0) = \gamma(a) $$
    and similar for $b$, so $\Gamma$ is proper.
    \qed

\eproof

\blemm

    Let $\Gamma\colon J\times[a,b]\to M$ be an admissible family of curves.
    Then on every rectangle $J\times[a_i,a_{i+1}]$ where $\Gamma$ is smooth,
    $$ D_s\partial_t\Gamma = D_t\partial_s\Gamma . $$

\elemm

\bproof

    This is local so we compute in local coordinates $(x^i)$ around a point $\Gamma(s_0,t_0)$.
    Writing
    $$ \Gamma(s,t) = (x^1(s,t),\dots,x^n(s,t)) , $$
    we have
    $$ \partial_t\Gamma = \pdvof{x^k}t\partial_k,\qquad \partial_s\Gamma = \pdvof{x^k}s\partial_k . $$
    Then using the formula for the covariant derivative,
    $$ \eqalign{
        D_s\partial_t\Gamma &= \parens{\pdvof{^2x^k}{st}+\pdvof{x^i}t\pdvof{x^j}s\Gamma^k_{ji}}\partial_k,\cr
        D_t\partial_s\Gamma &= \parens{\pdvof{^2x^k}{ts}+\pdvof{x^i}s\pdvof{x^j}t\Gamma^k_{ji}}\partial_k.
    } $$
    By symmetry of the second partial derivative and the connection, we see that these are equal.
    \qed

\eproof

\subsection{The curvature tensor}

Let $\bar\nabla$ be the standard connection on $\bR^n$, then a direct computation shows
$$ \bar\nabla_X\bar\nabla_YZ - \bar\nabla_Y\bar\nabla_XZ = \bar\nabla_{[X,Y]}Z . $$
This can be shown to be tensorial (which we will do), and thus we need only verify this for coordinate frames.
And since $\bar\nabla_{\partial_i}\partial_j=0$ this is clear.

This means that any Riemannian manifold which is flat, meaning locally isometric to $\bR^n$, must satisfy the above
identity with its own Levi-Civita connection.

\bdefn

    Let $(M,g)$ be Riemannian, then define a map $R\colon\X(M)\times\X(M)\times\X(M)\to\X(M)$ by
    $$ R(X,Y)Z = \nabla_X\nabla_YZ - \nabla_Y\nabla_XZ - \nabla_{[X,Y]}Z . $$

\edefn

\bprop

    $R$ is multilinear over $C^\infty(M)$ and thus defines a $(1,3)$-tensor field on $M$.

\eprop

\bproof

    Linearity in $Z$ is clear, and multilinearity over $\bR$ is as well.
    We show linearity in $Y$:
    $$ \eqalign{
        R(X,fY)Z &= \nabla_X\nabla_{fY}Z - \nabla_{fY}\nabla_XZ - \nabla_{[X,fY]}Z\cr
            &= \nabla_X(f\nabla_YZ) - f\nabla_Y\nabla_XZ - \nabla_{f[X,Y]+(Xf)Y}Z\cr
            &= (Xf)\nabla_YZ + f\nabla_X\nabla_YZ - f\nabla_Y\nabla_XZ - f\nabla_{[X,Y]}Z - (Xf)\nabla_YZ\cr
            &= fR(X,Y)Z .
    } $$
    Linearity in $X$ is done similarly.
    \qed

\eproof

We adopt the convention that the last index in $R$ is its contravariant (upper) one.
So in terms of a coordinate frame,
$$ R = R_{ijk}^\ell dx^i\otimes dx^j\otimes dx^k\otimes\partial_\ell,\qquad
R(\partial_i,\partial_j)\partial_k = R_{ijk}^\ell\partial_\ell . $$

\bprop

    The components of $R$ in any coordinate frame are given by
    $$ R_{ijk}^\ell = \partial_i\Gamma^\ell_{jk} - \partial_j\Gamma^\ell_{ik} + \Gamma^m_{jk}\Gamma^\ell_{im}
    - \Gamma^m_{ik}\Gamma^\ell_{jm} . $$

\eprop

\bproof

    A standard computation.
    \qed

\eproof

\bnote

    Curvature in the one-dimensional case is uninteresting.
    This is because $R(X,Y)=-R(Y,X)$, and in the one-dimensional case this clearly implies $R=0$.
    In particular if we consider the image of a curve to be a one-dimensional manifold (when the curve is not self-intersecting), then its curvature
    is zero.

\enote

\bprop

    Let $\Gamma$ be a smooth one-parameter family of curves $J\times I\to M$.
    Then for every vector field $V$ along $\Gamma$,
    $$ D_sD_tV - D_tD_sV = R(\partial_s\Gamma,\partial_t\Gamma)V . $$

\eprop

\bproof

    Locally write
    $$ \Gamma(s,t) = (\gamma^1(s,t),\dots,\gamma^n(s,t)),\qquad V(s,t) = V^j(s,t)\partial_j|_{\Gamma(s,t)} . $$
    Then
    $$ D_tV = \pdvof{V^i}t\partial_i + V^iD_t\partial_i , $$
    and so
    $$ D_sD_tV = \pdvof{^2V^i}{st}\partial_i + \pdvof{V^i}tD_s\partial_i + \pdvof{V^i}sD_t\partial_i + V^iD_sD_t\partial_i . $$
    Thus we get
    $$ D_sD_tV - D_tD_sV = V^i(D_sD_t\partial_i - D_tD_s\partial_i) . $$

    Now we need to compute the commutator.
    Let us write
    $$ S = \partial_s\Gamma = \pdvof{\gamma^k}s\partial_k,\qquad T = \partial_t\Gamma = \pdvof{\gamma^j}t\partial_j . $$
    Because $\partial_i$ is extendible,
    $$ D_t\partial_i = \nabla_T\partial_i = \pdvof{\gamma^j}{\partial t}\nabla_{\partial_j}\partial_i . $$
    Since $\nabla_{\partial_j}\partial_i$ is also extendible,
    $$ \eqalign{
        D_sD_t\partial_i &= 
            \pdvof{^2\gamma^j}{st}\nabla_{\partial_j}\partial_i + \pdvof{\gamma^j}t\nabla_S(\nabla_{\partial_j}\partial_i)\cr
            &= \pdvof{^2\gamma^j}{st}\nabla_{\partial_j}\partial_i + \pdvof{\gamma^j}t\pdvof{\gamma^k}s\nabla_{\partial_k}
                \nabla_{\partial_j}\partial_i .
    } $$
    We then see that
    $$ D_sD_t\partial_i - D_tD_s\partial_i =
    \pdvof{\gamma^j}t\pdvof{\gamma^k}s(\nabla_{\partial_k}\nabla_{\partial_j}\partial_i -
        \nabla_{\partial_j}\nabla_{\partial_k}\partial_i) = \pdvof{\gamma^j}t\pdvof{\gamma^k}sR(\partial_k,\partial_j)
        = R(S,T)\partial_i . $$
    Thus
    $$ D_sD_tV - D_tD_sV = V^iR(S,T)\partial_i = R(S,T)V $$
    as desired.
    \qed

\eproof

\bdefn

    Let $(M,g)$ be a Riemannian manifold.
    Define its {\emphcolor Riemannian curvature tensor} to be the tensor obtained by lowering $R$'s index:
    $\Rm=R^\flat$.
    Explicitly, it is a $(0,4)$-tensor field given by
    $$ \Rm(X,Y,Z,W) = \iprod{R(X,Y)Z,W} . $$

\edefn

In terms of local coordinates,
$$ \Rm = R_{ijk\ell}dx^i\otimes dx^j\otimes dx^k\otimes dx^\ell,\qquad R_{ijk\ell} = g_{\ell m}R_{ijk}^m . $$
Thus
$$ R_{ijk\ell} = g_{\ell m}(\partial_i\Gamma^m_{jk} - \partial_j\Gamma^m_{ik} + \Gamma^p_{jk}\Gamma^m_{ip} -
\Gamma^p_{ik}\Gamma^m_{jp}) . $$

Riemannian curvature is preserved by local isometries.

\bprop

    If $\phi\colon M\to\tilde M$ is a local isometry, let $\Rm,\tilde\Rm$ be the Riemannian curvature tensors of $M,\tilde M$
    respectively.
    Then $\phi^*\tilde\Rm=\Rm$.

\eprop

\bproof

    We have
    $$ \phi^*\tilde\Rm(X,Y,Z,W)_p = \tilde\Rm_{\phi p}(d\phi_pX,d\phi_pY,d\phi_pZ,d\phi_pW)
    = \iprod{R_{\phi p}(d\phi_pX,d\phi_pY)d\phi_pZ,d\phi_pW} . $$
    Then using the definition of the curvature tensor and the naturality of Levi-Civita, we get the desired result.
    \qed

\eproof

\bprop

    The following identities hold for the Riemannian curvature tensor:
    \benum
        \item $\Rm(W,X,Y,Z)=-\Rm(X,W,Y,Z)$,
        \item $\Rm(W,X,Y,Z)=-\Rm(W,X,Z,Y)$,
        \item $\Rm(W,X,Y,Z)=\Rm(Y,Z,W,X)$,
        \item $\Rm(W,X,Y,Z)+\Rm(X,Y,W,Z)+\Rm(Y,W,X,Z)=0$.
    \eenum

\eprop

\bproof

    The first is immediate because $R(W,X)Y=-R(X,W)Y$.
    To prove $Y$ it is sufficient to show that $\Rm(W,X,Y,Y)=0$ for all $Y$ (equivalence of antisymmetric multilinear
    maps).
    This follows from expanding $WX\norm Y^2,XW\norm Y^2,[W,X]\norm Y^2$ and equating them.

    To prove the fourth, it is sufficient to show
    $$ R(W,X)Y + R(X,Y)W + R(Y,W)X = 0 , $$
    which can be shown via a tedious computation.

    Then the third follows from the others, by a tedious expansion.
    \qed

\eproof

\subsection{Minimizing curves}

\bdefn

    If $I,K\subseteq\bR$ are intervals, $\gamma\colon I\to M$ a geodesic, and $\Gamma\colon K\times I\to M$ a variation of $\gamma$.
    It is called a {\emphcolor variation through geodesics} if its main curves $\Gamma_s(t)=\Gamma(s,t)$ are all geodesics
    (requiring $\Gamma$ be smooth).

\edefn

\bthrm[title=Jacobi equation]

    Let $\gamma$ be a geodesic in $M$ and $J$ a vector field along $\gamma$.
    If $J$ is the variation field of a variation through geodesics, then $J$ satisfies the {\emphcolor Jacobi equation}:
    $$ D_t^2J + R(J,\gamma')\gamma' = 0 . $$

\ethrm

\bproof

    Let $\Gamma$ be a variation through geodesics of $\gamma$ such that $J$ is its variation field, meaning $J(t)=S(0,t)$.
    Let $T=\partial_t\Gamma(s,t)$ and $S=\partial_s\Gamma(s,t)$ so $T$ is a geodesic.
    Then $D_tT=0$ identically.

    Then if we take the covariant derivative with respect to $s$ we get $D_sD_tT=0$.
    Recall curvature measures the non-commutativity of covariant derivatives, so
    $$ 0 = D_sD_tT = D_tD_sT + R(S,T)T = D_tD_tS + R(S,T)T , $$
    where the last equality is due to symmetry.
    Evaluating at $s$ gives us
    $$ D_t^2J + R(J,\gamma')\gamma' $$
    since $\gamma'=T(0,t)$.
    \qed

\eproof

\bdefn

    A smooth vector field along a geodesic satisfying the Jacobi equation is called a {\emphcolor Jacobi field}.

\edefn

\blemm

    If $\gamma\colon I\to M$ is geodesic, starting at $v$ and $J\in\X(\gamma)$ a Jacobi field, then
    $$ \iprod{\gamma',J} = \iprod{v,J(0)} + \iprod{v,J'(0)}t . $$

\elemm

\bproof

    Differentiating
    $$ \drv t\iprod{\gamma',J} = \iprod{D_t\gamma',J} + \iprod{\gamma',D_tJ} = \iprod{\gamma',J'} . $$
    And
    $$ \drv t\iprod{\gamma',J'} = \iprod{\gamma',J''} = - \iprod{\gamma',R(\gamma',J)\gamma'} , $$
    where last equality is due to Jacobi equation.
    $$ = \Rm(\gamma',J,\gamma',\gamma') = 0 , $$
    by symmetry of the Riemannian curvature tensor.
    Thus $\iprod{\gamma',J}=a+bt$.
    Evaluating the function and its derivative at $0$ gives us
    $$ a = \iprod{\gamma'(0),J(0)},\qquad b = \iprod{\gamma'(0),J'(0)} , $$
    as desired.
    \qed

\eproof

\blemm[title=Gauss]

    Let $p\in(M,g)$, and $w,v\in T_pM$, then
    $$ \iprod{d(\exp_p)_v(v),d(\exp_p)_v(w)}_{\exp_p(v)} = \iprod{v,w}_p . $$
    Where $d(\exp_p)_v\colon T_vT_pM\cong T_pM\to T_{\exp_p(v)}M$ under the natural identification $T_vT_pM\cong T_pM$.

\elemm

\bproof

    Define
    $$ f\colon(-\epsilon,\epsilon)\times[0,1] \to M,\qquad (s,t)\mapsto\exp_p(t(v+sw)) . $$
    So $f(s,1)=\exp_p(v+sw)$ and $t\mapsto f(t,s)$ is a geodesic for every $s$, and thus it is a variation through geodesics.
    Thus $J=\partial_sf(0,t)$ is a Jacobi field along $\gamma(t)=\exp_p(tv)$.
    Then
    $$ J(1) = \pdvof fs(0,1) = \drv s\exp_p(v+sw)|_{s=0} = d(\exp_p)_v(w) , $$
    and
    $$ \gamma'(1) = d(\exp_p)_v(v) . $$
    Thus by the above lemma
    $$ \iprod{d(\exp_p)_v(v),d(\exp_p)_v(w)} = \iprod{\gamma',J}(1) = \iprod{v,J(0)} + \iprod{v,J'(0)} . $$
    So we compute $J,J'$ at zero.
    $$ J(0) = \pdvof fs(0,0) = \drv s\exp_p(0)|_{s=0} = 0 , $$
    since $\exp_p(0)=p$ is constant.
    And
    $$ J'(0) = \frac D{\partial t}\pdvof fs(0,0) = \eval{\frac D{\partial s}}_{s=0}(\evpdv t_{t=0}f) . $$
    $\exp_p(t(v+sw))$ is a geodesic starting at $v+sw$, so its derivative with respect to $t$ at zero is $v+sw$, so this
    is equal to
    $$ J'(0) = \eval{\frac D{\partial s}}_{s=0}(v+sw) = w . $$
    Thus
    $$ \iprod{d(\exp_p)_v(v),d(\exp_p)_v(w)} = \iprod{v,w} $$
    as desired.
    \qed

\eproof

\bthrm

    Let $B\subseteq M$ be a geodesic ball around $p\in M$, and let $\gamma\colon I\to B$ be a geodesic starting at $p$.
    Let $c\colon I\to M$ be a piecewise smooth curve such that $c(0)=p,c(1)=\gamma(1)$.

    Then $L_g(\gamma)\leq L_g(c)$, and if there is equality then $c$ is a reparameterization of $\gamma$, i.e.\ there is a
    non-decreasing $\phi\colon I\to I$ such that $c=\gamma\circ\phi$.

\ethrm

\bproof

    We assume $c(t)\in B$ for all $t$, and $c(t)\neq p$ for all $t\in(0,1]$.
    If $c(t)=p$, then we can consider $c|_{[t,1]}$.
    Then we have since $c(t)\in B$ and $\exp_p\colon B_\epsilon(0)\to B$ is a diffeomorphism,
    $$ c(t) = \exp_p(w(t)),\qquad w(t)\in T_pM $$
    for a smooth $w\colon I\to B_\epsilon(0)$.
    Since $c(t)\neq p$ we have $w(t)\neq0$ for $t>0$.
    Then define smooth maps
    $$ r(t) = \norm{w(t)},\qquad v(t) = \frac{w(t)}{r(t)} . $$
    Now define
    $$ f\colon[0,\epsilon)\times[0,1]\to M,\qquad f(r,t) = \exp_p(rv(t)) . $$
    Then
    $$ c(t) = f(r(t),t) , $$
    so
    $$ c'(t) = \pdvof fr\cdot r' + \pdvof ft , $$
    and thus
    $$ \norm{c'(t)}^2 = \norm{\pdvof fr}^2\abs{r'}^2 + \norm{\pdvof ft}^2 + 2r'\iprod{\pdvof fr,\pdvof ft} . $$
    Notice that
    $$ \eqalign{
        \iprod{\pdvof fr,\pdvof ft} &= \iprod{\drv r\exp(rv(t)),\drv t\exp(rv(t))}\cr
            &= \iprod{d(\exp)_{rv(t)}(v(t)),d(\exp)_{rv(t)}(rv'(t))}\cr
            &= r\iprod{v(t),v'(t)}_p
    } $$
    where the last equality is due to Gauss' lemma.
    This is equal to
    $$ = \frac r2\drv t\iprod{v(t),v(t)} = 0 , $$
    since $\norm{v(t)}=0$ and thus its derivative is zero.

    Thus we have
    $$ \norm{c'(t)}^2 = \norm{\pdvof fr}^2\abs{r'}^2 + \norm{\pdvof ft}^2 \geq \norm{\pdvof fr}^2\abs{r'}^2 . $$
    Furthermore, let $\beta(r)=\exp(rv(t))$ so it is a geodesic with initial velocity $v(t)$, and
    $$ \pdvof fr = \evpdv r_{(r,t)}\exp(rv(t)) = \beta'(r) . $$
    Since geodesics have constant speed,
    $$ \norm{\pdvof fr} = \norm{\beta'(r)} = \norm{\beta'(0)} = \norm{v(t)} = 1 . $$
    Thus
    $$ \norm{c'(t)} \geq \abs{r'} . $$

    Furthermore we have
    $$ \gamma(1) = c(1) = \exp_p(r(1)v(1)) , $$
    and since $\gamma$ is a geodesic we must then have
    $$ \gamma(t) = \exp_p(tr(1)v(1)) \implies \norm{\gamma'(t)} = \norm{\gamma'(0)} = r(1)v(1) . $$
    Thus we have
    $$ L_g(\gamma) = \int_0^1\norm{r(1)v(1)} = \abs{r(1)} = r(1) . $$
    And (recall that the following is an improper integral)
    $$ L_g(c) = \int_0^1\norm{c'(t)} \geq \int_0^1 \abs{r'} \geq \int_0^1 r' = r(1) - \lim_{t\to0}r(t) = r(1) . $$
    So we have the desired inequality.

    If we have equality then we must have that $\norm{\pdvof ft}=0$ so $\pdvof ft=0$.
    That is
    $$ 0 = \pdvof ft = d\exp_p(rv'(t)) $$
    which implies, because $\exp_p$ is a diffeomorphism, that $v'(t)=0$.
    Thus $v(t)=v(1)$ is constant.
    Furthermore, we must have (by considering our inequalities above) that $\abs{r'}=r'$ so we must have $r'\geq0$.
    So
    $$ c(t) = \exp_p(r(t)v(1)) . $$
    And since $\gamma(t)=\exp_p(tr(1)v(1))$ this implies $c(t)=\gamma(r(t)/r(1))$ and $r(t)/r(1)$ has nonnegative
    derivative and thus is nondecreasing, as desired.

    Finally in the case that $c$ leaves the ball $B$, let $t_1\in[0,1]$ be the first value such that $c(t)\in\cl B-B$
    (by compactness one exists).
    Thus $L_g(c|_{[0,t_1)})$ is greater than the length of the radial geodesic from $p$ to $c(t_1)$ by our above
    proof.
    This has length equal to the radius of $B$ (as it is radial), and since $\gamma$ lies within $B$,
    its length is less than the radius.
    \qed

\eproof

Notice that this gives another proof of the Riemannian distance being a metric.
Because if $p\neq q$ then any curve connection $p$ and $q$ has curve length greater than a geodesic defined in a geodesic ball around $p$,
which is nonzero.

Recall that the topology on $M$ is equivalent to the metric topology induced by the metric $d_g$.
Thus $d_g$ is a continuous map $M\times M\to\bR$.

\bcoro

    If $B$ is a geodesic ball around $p$, then for every $q\in B$ the unique minimizing curve (up to reparameterization) from $p$ to $q$ is given by
    the radial geodesic from $p$ to $q$.

\ecoro

\bnote

    Since $q\in B$ we have that $q=\exp_p(w)$ for some $w\in B_\epsilon(0)$ (the ball mapped to $B$ diffeomorphically by $\exp_p$).
    The {\emphcolor radial geodesic} from $p$ to $q$ is the geodesic $\gamma\colon[0,1]\to M$ given by $t\mapsto\exp_p(tw)$.

\enote


\bdefn

    A {\emphcolor uniformly normal neighborhood} of $p\in M$ is an open neighborhood $p\in W$ such that there is a $\delta>0$ such that for
    all $q\in W$, $W\subseteq\exp_q(B_\delta(0))$.

\edefn

\blemm

    Every point in $M$ has a uniformly normal neighborhood.

\elemm

\bproof

    For $p\in M$, an open neighborhood $p\in V$, and $\epsilon>0$ define
    $$ U_{V,\epsilon} = \set{(q,v)}[q\in V,v\in T_qM,\norm v<\epsilon] \subseteq TM|_V . $$
    For small enough $\epsilon$, $\exp$ is defined on all of $U_{V,\epsilon}$.
    So the following map is smooth
    $$ F\colon U_{V,\epsilon} \to M\times M,\qquad (q,v)\mapsto(q,\exp_q(v)) . $$
    Let $\phi\colon U\to M$ be a coordinate chart which we identify with its coordinate chart $\phi\colon U\times\bR^n\to TM$
    of $TM$:
    $$ \phi(u,v) = (\phi(u),v^i\partial_i|_u) . $$
    Then
    $$ F\circ\phi(u,v) = (\phi(u),\exp_{\phi(u)}(v^i\partial_i|_u)) . $$
    Thus
    $$ d(F\circ\phi)_0 = \pmatrix{d\phi_0 & 0\cr d\phi_0 & d\phi_0} , $$
    which is invertible, and thus by the inverse function theorem there is a neighborhood of $(p,0)$ such that $F$ is a diffeomorphism
    onto its image.

    Since $U_{V,\epsilon}$ form a local basis around $(p,0)\in TM$, we can choose this neighborhood to be $U_{V,\delta}$.
    So $F(U_{V,\delta})$ is open in $M\times M$, there is an open neighborhood $p\in W$ such that $W\times W\subseteq F(U_{V,\delta})$.
    So for every $q\in W$ we have that
    $$ \set q\times W\subseteq F(U_{V,\delta})\cap(\set q\times M) = \set q\times\exp_q(B_\delta(0)) . $$
    Thus $W\subseteq\exp_q(B_\delta(0))$ as desired.
    \qed

\eproof

\bdefn

    An admissible curve $\gamma$ is {\emphcolor minimizing} if $L_g(\gamma)\leq L_g(\tilde\gamma)$ for every other admissible curve $\tilde\gamma$
    with the same endpoints.

    And $\gamma\colon I\to M$ is {\emphcolor locally minimizing} if for every $t_0\in I$ there is a neighborhood $t_0\in I_0\subseteq I$
    such that whenever $a<b$ are in $I_0$ then $\gamma|_{[a,b]}$ is minimizing.

\edefn

Note that $\gamma\colon[a,b]\to M$ is minimizing iff $d_g(\gamma(a),\gamma(b))=L_g(\gamma)$ by definition of $d_g$.

\blemm

    Let $p\in M$ and $B$ be a geodesic ball around $p$.
    Then for every $q\in B$, the radial geodesic from $p$ to $q$ is the unique minimizing curve from $p$ to $q$.

\elemm

\bproof

    Suppose $q=\exp_p(cv)$ for $v\in T_pM$ a unit vector.
    Then $\gamma\colon[0,c]\to M$ given by $\gamma(t)=\exp_p(tv)$ is the radial geodesic from $p$ to $q$ and it has arc length $c$.

    Now let $\sigma\colon[0,b]\to M$ be admissible from $p$ to $q$ which we can assume is parameterized by arc length.
    Let $a_0\in[0,b]$ be the last time $\sigma(t)=p$, and $b_0\in[0,b]$ the first time after $a_0$ that $\sigma(t)$ meets the geodesic
    sphere $S_c$ of radius $c$ around $p$.
    Then 

\eproof

\bthrm

    Every geodesic is locally minimizing.
    Conversely, every admissible minimizing curve is a geodesic (in particular it is smooth).

\ethrm

\bproof

    Let $\gamma\colon I\to M$ be a geodesic, which we take to be defined on an open interval.
    Let $t_0\in I$ and let $W$ be a uniformly normal neighborhood of $\gamma(t_0)$.
    And let $I_0\subseteq I$ be the connected component of $\gamma^{-1}W$ containing $t_0$.
    If $a<b$ are in $I_0$ then the definition of a uniformly normal neighborhood implies $\gamma(b)$ is contained in a geodesic ball
    centered at $\gamma(a)$.
    Thus there is a unique minimizing curve from $\gamma(a)$ to $\gamma(b)$ given by a radial geodesic.
    Since $\gamma|_{[a,b]}$ is also a geodesic segment from $\gamma(a)$ to $\gamma(b)$ lying in the same geodesic ball, it must coincide
    with this radial geodesic.
    So $\gamma|_{[a,b]}$ is minimizing for all $a<b$ in $I_0$, and thus $\gamma$ is locally minimizing.

    Conversely suppose $\gamma\colon I=[a,b]\to M$ is a minimizing admissible curve.
    Then for every $t_0\in I$ as before we can find a connected neighborhood $I_0$ of $t_0$ such that $\gamma(I_0)$ is contained
    in a uniformly normal neighborhood $W$.
    Then for every $a_0,b_0\in I_0$, the unique minimizing curve $\gamma(a_0)$ to $\gamma(b_0)$ is a radial geodesic.
    Since $\gamma|_{[a_0,b_0]}$ is minimizing, it must coincide with this radial geodesic.
    So $\gamma$ is a geodesic in a neighborhood of $t_0$ for all $t_0\in I$, since being a geodesic is a local property, $\gamma$ is a geodesic.
    \qed

\eproof

Note that there does not need to exist an admissible minimizing curve between two points, even if the manifold is connected.

\subsection{Sectional curvature}

\blemm

    Let $p\in M$ and $\sigma\subseteq T_pM$ be a two-dimensional subspace.
    Then for every basis $x,y$ of $\sigma$, define
    $$ K(x,y) = \frac{\Rm(x,y,x,y)}{\iprod{x,x}\iprod{y,y}-\iprod{x,y}^2} . $$
    This is a constant, independent of choice of basis.

\elemm

\bproof

    Every basis of $\sigma$ can be obtained from $x,y$ by elementary operations:
    $$ x,y\mapsto y,x,\qquad x,y\mapsto x,y\mapsto x,\lambda y,\qquad x,y\mapsto x+\lambda y,y . $$
    Since $\Rm$ is antisymmetric in both its first two and last two coordinates, $K$ is invariant under the first elementary operation
    (swapping).
    And by multilinearity it is routine to check that it is invariant under the second two elementary operations.
    Thus $K$ is constant.
    \qed

\eproof

Note that $\iprod{x,x}\iprod{y,y}-\iprod{x,y}^2$ is the square of the area of parallelogram spanned by $x,y$ and is nonzero when $x,y$ are
linearly independent.

\bdefn

    For a Riemannian manifold $(M,g)$, we define its {\emphcolor sectional curvature} at $(p,\sigma)$ for $\sigma\subseteq T_pM$
    a two dimensional subspace by
    $$ K_p(\sigma) = K(x,y)\qquad\hbox{for any basis $x,y\in\sigma$} . $$

\edefn

\blemm

    Let $V$ be a finite-dimensional vector space and $F,G\colon V^4\to\bR$ multilinear functionals.
    Further suppose that $F,G$ satisfy the symmetries of Riemannian curvature.
    If $F(x,y,x,y)=G(x,y,x,y)$ for all $x,y\in V$ then $F=G$.

\elemm

\bproof

    Substitute $y+z$ for $y$ in $F(x,y,x,y)=G(x,y,x,y)$ and use the symmetries to obtain
    $$ F(x,y,x,z) = G(x,y,x,z) . $$
    Then substitute $x+w$ for $x$ to obtain
    $$ F(x,y,w,z) = G(x,y,w,z) . \qed $$

\eproof

In particular this means that sectional curvature determines Riemannian curvature.
That is, if we are given a Riemannian manifold, it is enough to know either the Riemannian or sectional curvature to obtain the other.

\bcoro

    Let $M$ be a Riemannian manifold, then define a $4$-tensor on $M$ by
    $$ \Rm'(X,Y,Z,W) = \iprod{X,W}\iprod{Y,Z} - \iprod{Y,W}\iprod{X,Z} . $$
    Then $M$ has constant sectional curvature $K_0$ at $p$ iff $\Rm=K_0\Rm'$ at $p$.

\ecoro

\bproof

    Note that $\Rm'$ satisfies the symmetries of Riemannian curvature, then $\Rm(X,Y,X,Y)=K_0\Rm'(X,Y,X,Y)$ by definition.
    By the above lemma this means $\Rm=K_0\Rm'$ as desired, and the converse is immediate.
    \qed

\eproof

\subsection{Jacobi fields}

Recall that a {\emphcolor Jacobi field} is a vector field $J$ along a geodesic $\gamma$ satisfying the {\emphcolor Jacobi equation}:
$$ D_t^2J + R(J,\gamma')\gamma' = 0 . $$

\bprop

    Let $\gamma\colon I\to M$ be a geodesic on a Riemannian manifold $M$.
    Let $t_0\in I$ and $p=\gamma(t_0)$, then for every $v,w\in T_pM$ there is a unique Jacobi field $J$ along $\gamma$ such that
    $$ J(a) = v,\qquad D_tJ(a) = w . $$

\eprop

\bproof

    Take a parallel orthonormal frame $(E_i)$ along $\gamma$ (obtained by parallel transport of an orthonormal basis), and write
    $v=v^iE_i(a)$ and $w=w^iE_i(a)$, and $\gamma'=y^iE_i$.
    Then $J\in\X(\gamma)$ can be written as $J=J^jE_j$, so that $D_t^2J=\ddot J^jE_j$ since the $E_j$ are paralell.
    The Jacobi equation becomes
    $$ 0 = \ddot J^iE_i + R(J^jE_j,y^kE_k)(y^\ell E_\ell) = \ddot J^iE_i + J^jy^ky^\ell R_{jk\ell}^i(\gamma(t))E_i , $$
    where $R_{jk\ell}^i(\gamma)$ are the components of $R$ in the frame $E_i$.
    This is zero iff for all $i$,
    $$ 0 = \ddot J^i + J^jy^ky^\ell R_{jk\ell}^i(\gamma) , $$
    Notice that this is all smooth, and thus a smooth second-order linear ODE, which has a unique solution along all of $I$.
    \qed

\eproof

Notice that for a geodesic $\gamma$, $\dot\gamma$ itself is a Jacobi field along $\gamma$.
Indeed, $\dot\gamma$ is the variation field of $\Gamma(s,t)=\gamma(s)$, and variation fields of variations through geodesics are Jacobi.

Further note that the collection of Jacobi fields along $\gamma$ form a subspace of $\X(\gamma)$, by the linearity of the Jacobi field
equation.

\bcoro

    Let $(M,g)$ be a Riemannian manifold of dimension $n$, and $\gamma$ a geodesic on $M$.
    The set of Jacobi fields along $\gamma$ form a $2n$-dimensional subspace of $\X(\gamma)$.

\ecoro

\bproof

    For any $a\in I$, let $p=\gamma(a)$, and consider the map from Jacobi fields along $\gamma$ to $T_pM\oplus T_pM$ given by
    $J\mapsto(J(a),D_tJ(a))$.
    The above proposition shows that this is a bijection.
    \qed

\eproof

We consider the case of a Riemannian manifold with constant sectional curvature.
For $K\in\bR$ let $s_K\colon\bR\to\bR$ be defined by
$$ s_K(t) = \cases{t&if $K=0$\cr\noalign{\kern2pt} R\sin t/R&if $K=1/R^2>0$\cr\noalign{\kern2pt} R\sinh t/R&if $K=-1/R^2<0$.} $$

\blemm

    Suppose $(M,g)$ is a Riemannian manifold with constant sectional curvature $K$.
    Let $\gamma$ be a unit-speed geodesic in $M$.
    Then the Jacobi fields $J$ orthogonal to $\dot\gamma$ which vanish at $t=0$ are given by
    $$ J(t) = c\cdot s_K(t)E(t) , $$
    where $c$ is an arbitrary constant, and $E\in\X(\gamma)$ is any vector field parallel along $\gamma$ and orthogonal to $\dot\gamma$.

\elemm

\bproof

    Since $M$ has constant curvature, we have that
    $$ \Rm(X,Y,Z,W) = K(\iprod{X,W}\iprod{Y,Z} - \iprod{Y,W}\iprod{X,Z}) . $$
    So we have
    $$ R(X,Y)Z = K(\iprod{Y,Z}X - \iprod{X,Z}Y) . $$
    Substituting into the Jacobi equation we get
    $$ 0 = D_t^2J + K(\iprod{\gamma',\gamma'}J - \iprod{J,\gamma'}\gamma') = D_t^2J + KJ . $$
    Choosing $J(t)=u(t)E(t)$, the ODE becomes
    $$ u''(t) + Ku(t) = 0 , $$
    whose solutions given $u(0)=0$ are multiples of $s_K$.
    Thus all fields of the form in the lemma are Jacobi fields.

    Now, the dimension of the space of Jacobi fields vanishing at $0$ orthogonal to $\gamma'$ is $n-1$ (the space of vanishing Jacobi fields has
    dimension $n$, by considering the isomorphism $J\mapsto(J(0),D_tJ(0))$, and we reduce the dimension by one by considering the orthogonal subspace
    to $\gamma'$).
    The space of vector fields parallel along $\gamma$ orthogonal to $\dot\gamma$ also has dimension $n-1$ (mapping $E$ to $E(0)\in T_pM$),
    so this gives all such Jacobi fields.
    \qed

\eproof

